\chapter{怎样真正把合数平均交给素数平均}\label{ch:transfer}
\section{为什么不能随时把六平均拆成二平均}
一个 $k$ 平均矩阵在 $n$ 维空间上的秩为 $n-k+1$。非空 $k$ 平均词的秩只会更小，所以不可能在\emph{整个空间}等于一个 $p<k$ 平均矩阵，其秩为 $n-p+1$。因此“六等于二乘三，所以能随时执行二平均”不成立。\src{S05}

但若每个商位置已有三份相同的值，两个商位置的六个实际副本恰好可以做一次六平均。这是在特殊子空间上模拟，完全可能。更一般地，$p\mid k$、每个商位置有 $c$ 份相同副本且 $k/p\mid c$ 时，一次商层 $p$ 平均用 $pc/k$ 次实际 $k$ 平均实现。

旧半块方法先制造这些相同副本，再调用二平均。奇数个商位置时，最后一半不能单独平均，便采用重叠入口：先平均 $2c$ 元组为 $a$，再取其中 $c$ 份与外部 $c$ 项平均为 $b=(a+v)/2$，留下 $c$ 份 $a$，形成商值 $(a,b,b,\ldots)$。在危险奇素数处，$a=b$ 等价于 $a=v$，所以见证被保留。\src{S05}

\section{后来的简化：起初块内不相等，也照样能模拟}
其实不必先让每块相等。把 $n$ 个位置分成等大的 $c$ 元块，只记录每块的\emph{均值} $z_i$。若选 $p$ 个块，对其全部 $pc$ 个位置执行一个固定全平均网络，它们的新均值就是
\[\frac1{pc}\sum_{\text{这些位置}}x_j=\frac1p\sum_{\text{这些块}}z_i.
\]
未选块均值不变。因此商均值严格跟随一条 $p$ 平均路径，被选块内部还顺便变得全等。\src{S06}

商均值最后都正确，还可能有某块一直没选过、内部不等。补一轮覆盖所有块的 $p$ 元块组即可：此时每个块均值相同，再平均不会改变目标，却会清掉所有块内差异。这一清理步骤不能省。

\section{先把危险素数从块均值里保护下来}
我们还须选择分块，使商输入合法。已有同时好划分引理说：块数 $m\ge5$，每个块大小在所有需保护素数处可逆，而且危险素数个数不超过 $m-1$，就能选到各素数处块均值不全相同的划分。\src{S29}

它的直观证明是逐个修复素数：交换计划中两个块的一个原位置，会改变两个块均值。每个已经保护的素数至多禁止一条块间边；删去少于 $m-1$ 条边的完全图仍连通。如果允许边两侧都找不到不同余位置，原输入就会全同余，矛盾。交换只是修改分组计划，最终仍对原位置实施平均。

\begin{theorem}[素数因子转译]\label{thm:factor}
取 $p\mid k$ 为素数，假设 $k/p\mid n$。定义
\[
c=p^{v_p(k)-1}\prod_{\ell\mid k,\ell\ne p}\ell^{v_\ell(n)},\qquad m=n/c.
\]
若 $m=p$，或 $p=2,m\ge4$，或 $p$ 奇且 $m\ge2p+1$，则 $\mathcal P(k,n)$ 成立。
\end{theorem}
\begin{proof}
$c\mid n$，$k\mid pc$，且 $pc$ 与 $k$ 有相同素因子支持。上一讲的张量网络可在任意 $pc$ 个位置实现全平均。
商维数 $m$ 的非 $p$ 素因子都不整除 $k$；它们是原输入已有的危险素数，并且 $c$ 在这些地方可逆。当 $m\ge5$ 时，用同时好划分安排 $m$ 个 $c$ 块；$m=4,p=2$ 时没有危险奇素数；$m=p$ 时直接对全部位置使用网络。
对合法商输入调用前篇素数定理，用 $pc$ 元网络逐步模拟其平均。所有块均值相同后，按本节开头的覆盖清理完成所有原位置。
\end{proof}

例如十五平均33元，取 $p=5,c=3,m=11$。每次选五个三元块，直接做一次十五平均；商层五平均十一元已解。十五平均35元则取 $p=3,c=5,m=7$，调用三平均七元。两者都是原来就有的位置，没有添加副本。

六平均14元取 $p=3,c=2,m=7$，六平均22元取 $p=3,c=2,m=11$，也都直接完成。这两项说明旧六平均未解表需要更新。

\section{哪些整除族已经全部解决}
\begin{corollary}
任意偶数 $k$ 的全部 $k$ 倍数维数均满足G判据。任意 $k\ge2$ 的全部 $n=ak$、$a\ge3$ 也满足G判据。
\end{corollary}
\begin{proof}
偶数取 $p=2$，商 $m$ 为偶数且至少二，故满足转译条件。
设 $k$ 奇，将倍数 $a=uv$ 分解为 $u$ 只含 $k$ 的素因子、$\gcd(v,k)=1$。对 $p\mid k$，商为 $m=p^{v_p(u)+1}v$。$v=1$ 时是 $p$ 或至少 $p^2$；$v\ge3$ 时至少 $3p$；$v=2$ 时因 $a\ge3$ 有 $u>1$，取 $p\mid u$，商至少 $2p^2$。均满足前篇阈值。
\end{proof}
更具体地，四平均覆盖全部偶数 $n\ge8$，另有 $n=4$；六平均覆盖全部三的倍数 $n\ge6$。后一结论中，去掉所有三因子后的商为二或至少四时用半块转译，若商为一，则用上一讲的 $3^b$ 固定网络。\src{S05}

这些都是整除族，不是连续区间。不能把“全部 $n=ak$、$a\ge3$”读成“全部 $n\ge3k$”。剩下不整除的邻点，正是接续问题。
